\chapter[2019 --- Jr. et al.]{2019: William G. Kaelin Jr., Sir Peter J. Ratcliffe, Gregg L. Semenza}
\label{ch:2019_Kaelin}

\section*{Main Contribution}

\textit{Discovered how cells sense and adapt to oxygen availability, explaining a fundamental adaptive mechanism critical for development, immunity, and cancer.}

\section{Official Citation}

for their discoveries of how cells sense and adapt to oxygen availability

\section{Background and Historical Context}

The 2019 Nobel Prize in Physiology or Medicine recognized William G. Kaelin Jr., Sir Peter J. Ratcliffe, and Gregg L. Semenza ``for their discoveries of how cells sense and adapt to oxygen availability.'' Their work illuminated one of the most fundamental adaptive mechanisms in biology---the way cells detect changes in oxygen levels and adjust their metabolic and physiological responses accordingly. This oxygen-sensing pathway plays essential roles in development, physiology, and disease, and its elucidation has opened up new therapeutic approaches for anemia, cancer, and ischemic diseases.

Oxygen is essential for life. Cells require a constant supply of oxygen to generate energy through aerobic metabolism, and even brief periods of oxygen deprivation (hypoxia) can cause irreversible cellular damage and death. Organisms have therefore evolved sophisticated mechanisms for detecting changes in oxygen levels and coordinating appropriate physiological responses, including changes in blood vessel growth (angiogenesis), red blood cell production (erythropoiesis), and metabolic reprogramming.

The importance of oxygen sensing was first recognized in the context of erythropoiesis---the production of red blood cells by the bone marrow. In the 1880s, it was observed that the hormone erythropoietin (EPO), produced primarily by the kidneys, was the primary regulator of red blood cell production. When oxygen levels in the blood decreased (hypoxia), the kidneys increased EPO production, which stimulated the bone marrow to produce more red blood cells, thereby increasing the oxygen-carrying capacity of the blood. When oxygen levels were restored to normal, EPO production decreased, and red blood cell production returned to baseline.

However, the molecular mechanism by which cells detected changes in oxygen levels and regulated EPO production remained unknown for more than a century. The identification of this oxygen-sensing mechanism became one of a major goals in physiology and molecular biology, with significant implications for understanding the regulation of erythropoiesis and the response to hypoxia.

\section{The Discovery/Contribution}

The elucidation of the oxygen-sensing pathway involved the combined efforts of Semenza, Ratcliffe, and Kaelin, working independently at different institutions but pursuing complementary lines of research.

\textbf{Gregg L. Semenza:} Semenza, an American physician-scientist born in 1955 in New York City, working at Johns Hopkins University, made the foundational discovery of the hypoxia-inducible factor (HIF). In the late 1980s and early 1990s, Semenza set out to identify the transcription factor responsible for the hypoxia-induced expression of the EPO gene. Using a combination of molecular biology and genetics, he identified a protein complex that bound to a specific DNA sequence---the hypoxia response element (HRE)---in the enhancer region of the EPO gene and activated its transcription under hypoxic conditions. He named this protein complex hypoxia-inducible factor (HIF).

Semenza showed that HIF consisted of two subunits: HIF-1$\alpha$, an oxygen-sensitive subunit that was rapidly degraded under normoxic (normal oxygen) conditions but stabilized under hypoxic conditions, and HIF-1$\beta$, a constitutively expressed subunit that was also known as the aryl hydrocarbon receptor nuclear translocator (ARNT). The oxygen-dependent regulation of HIF-1$\alpha$ was the key to the oxygen-sensing mechanism: under normoxic conditions, HIF-1$\alpha$ was continuously synthesized but rapidly degraded, while under hypoxic conditions, HIF-1$\alpha$ accumulated and formed active HIF-1 complexes with HIF-1$\beta$.

Semenza subsequently identified HIF-2$\alpha$, a second oxygen-sensitive subunit that was structurally similar to HIF-1$\alpha$ but had distinct tissue distributions and target gene specificities. While HIF-1$\alpha$ was ubiquitously expressed and regulated the expression of genes involved in glycolysis and other metabolic pathways, HIF-2$\alpha$ was more selectively expressed in specific cell types, including endothelial cells and renal interstitial fibroblasts, and played a particularly important role in the regulation of EPO production. The existence of two oxygen-sensitive HIF-$\alpha$ subunits with overlapping but distinct functions provided a mechanism for tissue-specific responses to hypoxia. HIF-1$\alpha$ and HIF-2$\alpha$ could heterodimerize with the same partner subunit (HIF-1$\beta$/ARNT) but activate different sets of target genes, allowing different tissues to mount tailored responses to hypoxic conditions.

Semenza also showed that HIF activated a broad program of gene expression that extended far beyond EPO. HIF target genes included those encoding vascular endothelial growth factor (VEGF), which promotes the growth of new blood vessels (angiogenesis); glycolytic enzymes, which enable cells to generate energy through glycolysis rather than oxidative phosphorylation; and glucose transporters, which increase glucose uptake into hypoxic cells. This coordinated transcriptional program enabled cells and tissues to adapt to reduced oxygen availability by promoting blood vessel growth, shifting to anaerobic metabolism, and enhancing glucose utilization.

\textbf{Sir Peter J. Ratcliffe:} Ratcliffe, a British physician-scientist born in 1954 in Lancashire, England, working at the University of Oxford and the Francis Crick Institute, made critical contributions to understanding the mechanism of HIF-1$\alpha$ degradation and the oxygen-sensing mechanism. Ratcliffe and his colleagues showed that HIF-1$\alpha$ was hydroxylated on specific proline residues under normoxic conditions, and that this hydroxylation was essential for its recognition and degradation by the von Hippel-Lindau (VHL) tumor suppressor protein.

The prolyl hydroxylase enzymes responsible for HIF-1$\alpha$ hydroxylation were identified by Ratcliffe's group and named hypoxia-inducible factor prolyl hydroxylases (HPHs) or prolyl hydroxylase domain proteins (PHDs). These enzymes used molecular oxygen as a co-substrate, making them direct sensors of cellular oxygen levels. Under normoxic conditions, when oxygen was abundant, the PHDs hydroxylated HIF-1$\alpha$, targeting it for ubiquitination by VHL and subsequent degradation by the proteasome. Under hypoxic conditions, when oxygen was scarce, PHD activity was reduced, allowing HIF-1$\alpha$ to accumulate and activate target gene transcription.

Ratcliffe also showed that HIF-1$\alpha$ was hydroxylated on an asparagine residue by the factor inhibiting HIF (FIH), which prevented the recruitment of the transcriptional coactivators p300 and CBP. This asparagine hydroxylation provided an additional layer of oxygen-dependent regulation of HIF activity, ensuring that HIF-mediated transcription was tightly controlled by oxygen levels. The identification of both prolyl and asparagine hydroxylation events established a two-tiered oxygen-sensing mechanism: prolyl hydroxylation controlled HIF-$\alpha$ protein stability (by targeting it for degradation), while asparagine hydroxylation controlled HIF transcriptional activity (by blocking coactivator recruitment). Both modifications depended on molecular oxygen as a co-substrate, making the PHDs and FIH direct sensors of cellular oxygen levels.

The elucidation of the hydroxylation mechanism also identified the 2-oxoglutarate-dependent dioxygenase superfamily as the enzyme family responsible for oxygen sensing. The PHDs and FIH belong to this superfamily, which also includes enzymes involved in collagen synthesis, fatty acid metabolism, and epigenetic regulation. The requirement for 2-oxoglutarate, iron, and ascorbate as co-factors for these enzymes provided additional regulatory inputs into the oxygen-sensing pathway and explained why conditions such as iron deficiency and scurvy could affect erythropoiesis and other hypoxia-responsive processes.

\textbf{William G. Kaelin Jr.:} Kaelin, an American physician-scientist born in 1957 in New York City, working at the Dana-Farber Cancer Institute and Harvard Medical School, made key contributions to understanding the role of the VHL tumor suppressor protein in HIF regulation. Kaelin had been studying the VHL gene, which is mutated in the hereditary cancer syndrome von Hippel-Lindau disease. Patients with VHL disease develop tumors in multiple organs, including clear cell renal cell carcinoma, hemangioblastomas, and pheochromocytomas.

Kaelin showed that the VHL protein was the substrate recognition component of an E3 ubiquitin ligase complex that targeted HIF-1$\alpha$ for proteasomal degradation. Mutations in VHL that disrupted its ability to bind HIF-1$\alpha$ led to constitutive stabilization of HIF-1$\alpha$ even under normoxic conditions, resulting in the activation of HIF target genes involved in angiogenesis (including vascular endothelial growth factor, VEGF), erythropoiesis (EPO), and metabolic reprogramming. This explained the highly vascular nature of VHL-associated tumors and provided a molecular link between the oxygen-sensing pathway and cancer biology.

The combined work of Semenza, Ratcliffe, and Kaelin revealed a comprehensive picture of the oxygen-sensing pathway. Under normoxic conditions, PHDs hydroxylate HIF-1$\alpha$ on proline residues, targeting it for VHL-mediated ubiquitination and proteasomal degradation. FIH hydroxylates HIF-1$\alpha$ on an asparagine residue, preventing transcriptional activation. Under hypoxic conditions, PHD and FIH activities are reduced due to limited oxygen availability, allowing HIF-1$\alpha$ to accumulate, form active complexes with HIF-1$\beta$ and the coactivators p300/CBP, and activate the transcription of target genes involved in erythropoiesis, angiogenesis, metabolic adaptation, and cell survival.

\section{Impact on Modern Medicine}

The elucidation of the oxygen-sensing pathway by Kaelin, Ratcliffe, and Semenza has had a significant impact on modern medicine, with applications in the treatment of anemia, cancer, and ischemic diseases.

\textbf{Anemia treatment:} The understanding of the HIF pathway has led to the development of a new class of drugs for the treatment of anemia associated with chronic kidney disease (CKD). HIF prolyl hydroxylase inhibitors (HIF-PHIs), also known as HIF stabilizers, are small molecules that inhibit the PHD enzymes, thereby preventing the degradation of HIF-1$\alpha$ and HIF-2$\alpha$ and stimulating the production of EPO. Roxadustat, daprodustat, and vadadustat are among the HIF-PHIs that have been approved or are in late-stage clinical trials for the treatment of CKD-associated anemia. These oral agents offer an alternative to injectable erythropoiesis-stimulating agents (ESAs) and may provide additional benefits through the pleiotropic effects of HIF activation, including improved iron metabolism and reduced inflammation.

\textbf{Cancer therapy:} The HIF pathway is frequently activated in cancer, particularly in hypoxic tumors, where it promotes angiogenesis, metabolic reprogramming, and therapy resistance. The understanding of the molecular mechanisms of HIF activation has identified new therapeutic targets for cancer, including HIF-1$\alpha$, HIF-2$\alpha$, and the PHD enzymes. HIF-2$\alpha$ inhibitors, including belzutifan (MK-6482), have shown promising results in clinical trials for VHL-associated renal cell carcinoma and other HIF-driven tumors. The FDA approved belzutifan in 2021 for the treatment of VHL-associated renal cell carcinoma, representing a direct clinical translation of the basic research on the oxygen-sensing pathway.

\textbf{Ischemic diseases:} The understanding of the HIF pathway has also informed the development of therapeutic strategies for ischemic diseases, including myocardial infarction, stroke, and peripheral artery disease. HIF activators, including PHD inhibitors, have been explored as potential treatments for ischemic conditions by promoting angiogenesis and improving tissue oxygenation. While these approaches are still in early stages of development, they represent promising strategies for the treatment of ischemic diseases.

\textbf{Wound healing:} The HIF pathway plays an important role in wound healing by promoting angiogenesis and tissue repair. HIF activators have been explored as potential treatments for chronic wounds, including diabetic foot ulcers and venous leg ulcers, which are characterized by impaired tissue oxygenation and defective healing. Early clinical results have been promising, suggesting that HIF-based therapies may have a role in the management of chronic wounds.

\textbf{Organ transplantation:} The understanding of the oxygen-sensing pathway has also informed strategies for improving organ preservation and transplantation. HIF-based approaches may improve the viability of donor organs by enhancing their tolerance to ischemia and reperfusion injury.

William G. Kaelin Jr., Sir Peter J. Ratcliffe, and Gregg L. Semenza were awarded the Nobel Prize in Physiology or Medicine in 2019 ``for their discoveries of how cells sense and adapt to oxygen availability.'' Their work has fundamentally changed our understanding of oxygen sensing and has opened up new therapeutic avenues for a range of diseases.

\section{Legacy}

The legacy of Kaelin, Ratcliffe, and Semenza's work extends across multiple areas of biology and medicine. Their discoveries have revealed one of the most fundamental adaptive mechanisms in biology---the ability of cells to sense and respond to changes in oxygen levels---and have provided a molecular framework for understanding the role of hypoxia in development, physiology, and disease.

The oxygen-sensing pathway is now recognized as a central regulator of cellular and organismal physiology, controlling not only erythropoiesis and angiogenesis but also metabolic reprogramming, stem cell function, immune regulation, and wound healing. The broad biological significance of the HIF pathway underscores the fundamental importance of the discoveries made by Kaelin, Ratcliffe, and Semenza.

The therapeutic potential of the oxygen-sensing pathway is still being explored, with new applications and new drugs in various stages of development. The approval of HIF-PHIs for anemia and HIF-2$\alpha$ inhibitors for cancer represents the beginning of a new era in the clinical application of the oxygen-sensing pathway.

The work of these three laureates also illustrates the power of interdisciplinary research and the importance of pursuing fundamental questions in biology. Their discoveries were made possible by the integration of genetics, biochemistry, molecular biology, and clinical medicine, and their impact has been felt across all of these fields.

Today, research on the oxygen-sensing pathway continues to yield new insights into the mechanisms of cellular adaptation and the pathogenesis of disease. The legacy of Kaelin, Ratcliffe, and Semenza's work provides the foundation for these ongoing investigations and will continue to shape the future of medicine and biology for generations to come.


\section{Modern Developments}

Kaelin, Ratcliffe, and Semenza's oxygen sensing work has led to modern cancer and anemia therapy. Understanding of HIF signaling has led to HIF-PHD inhibitors (roxadustat, daprodustat) for anemia. Understanding of HIF in cancer has led to research on HIF inhibitors for tumor metabolism. Understanding of oxygen sensing in immunity has revealed its role in macrophage function. Understanding of HIF in wound healing has informed development of hypoxia-inducible factor-based therapies.
